\chapter{Background Theory and Literature Review}
\label{ch:litreview}

This chapter surveys the literature underpinning the proposed pipeline and identifies the specific gap the dissertation addresses. It is condensed from a standalone literature review produced earlier in the project; the full nine-page version, including additional discussion of each publication, is available alongside this dissertation. The chapter is organised around eight thematic clusters (orbital data characterisation, super-resolution, image-to-image translation, cross-view synthesis, digital terrain modelling, neural scene rendering, VR reconstruction, and generative hallucination) before a comparative synthesis draws out the cross-cutting patterns that motivate the proposed architecture.

\section{Taxonomy and Conceptual Framework}

The literature reviewed here spans several independently advancing communities (remote sensing image processing, generative modelling, planetary science, and VR engineering) whose outputs are jointly necessary to the downstream synthesis requirement without any of them addressing it explicitly. Three axes organise the territory: \textit{spatial view} (nadir/orbital vs.\ perspective/rover vs.\ both), \textit{task type} (enhancement, translation, or geometric reconstruction), and \textit{output modality} (RGB texture, elevation/depth, or rendered video). Eight clusters emerge: orbital data characterisation~\cite{serrano2013hirise,osadnik2025mcted}; nadir super-resolution~\cite{tao2021cassis,wu2024rstsrn,lu2025thermal,khanna2024diffusionsat,zhu2025taming}; image-to-image translation~\cite{kim2024unsb,xiao2025brownian}; cross-view synthesis in Earth remote sensing~\cite{qian2023sat2density,li2024sat2scene,li2024crossviewdiff}; monocular digital terrain modelling~\cite{tao2021madnet,zou2026pfmgan,cao2024gen3d}; neural scene rendering~\cite{giusti2022marf,li2025martian,dealmeida2025neural3d,martinezpetersen2025hifi,paul2025nerfspace}; VR reconstruction~\cite{catalano2025inpainting,li2025martian}; and generative reliability~\cite{aithal2024hallucination,rathkopf2025reliability}.

The gap between the super-resolution cluster and the neural-rendering cluster is the primary concern of this dissertation. Nadir SR methods keep the viewpoint fixed; neural rendering methods require perspective imagery as input. No Mars-specific method reviewed here translates from nadir to perspective in a way that synthesises rover-viewable texture from orbital data. Section~\ref{sec:crossview} reports a cluster from Earth remote sensing that does address exactly this translation direction, absent from the dissertation's original reading list and added once this project's own CycleGAN results made clear it was needed.

\section{Super-Resolution for Planetary Remote Sensing}

Super-resolution research on planetary imagery has progressed through three generations since 2021. Tao \textit{et al.}~\cite{tao2021cassis} propose MARSGAN, a CaSSIS-to-HiRISE-equivalent enhancement of ESRGAN achieving roughly $3\times$ effective resolution improvement, evaluated only qualitatively across eight Mars sites. Co-registering CaSSIS and HiRISE precisely enough for pixel-level comparison remains non-trivial, so quantitative cross-method comparison stays opaque throughout this sub-field. Wu \textit{et al.}~\cite{wu2024rstsrn} apply recursive Swin Transformer attention (RSTSRN) at $\times2$--$\times8$, improving on LapSRN by up to 1.22\,dB PSNR at $\times8$, but evaluate only against bicubically downsampled inputs rather than the genuinely distinct point-spread function of a real orbital sensor, a protocol weakness that recurs across this literature and directly informed how the corpus for this dissertation's own CycleGAN baseline was later validated (Chapter~4).

Khanna \textit{et al.}~\cite{khanna2024diffusionsat} propose DiffusionSat, a diffusion foundation model for satellite imagery conditioned on geolocation and sensor metadata, demonstrating that multi-sensor domain generalisation is tractable within a diffusion framework, though its training corpus contains no Mars data and its conditioning assumes Earth coordinates. Zhu \textit{et al.}~\cite{zhu2025taming} address the fidelity/perceptual-realism trade-off directly relevant to preserving large-scale geomorphology while generating sub-metre texture, via frequency-decomposition conditioning that anchors denoising to low-frequency structure. Lu and Su~\cite{lu2025thermal} fuse THEMIS thermal imagery with visible and topographic data for SR, of indirect relevance since thermal is not the visible RGB domain required here.

Collectively, this literature shows $\times4$ enhancement is achievable with GAN, transformer, or diffusion backbones, but every reviewed method operates nadir-to-nadir: input and output share a viewpoint. The perspective transformation this dissertation requires is nowhere addressed by SR alone.

\section{Image-to-Image Translation as a Domain-Bridging Mechanism}

Image-to-image translation methods map entire image domains to each other, structurally better suited to a view transformation than SR. Kim \textit{et al.}~\cite{kim2024unsb} propose UNSB (Unpaired Neural Schr\"{o}dinger Bridge), reframing unpaired translation as learning an SDE that transports directly between two target distributions rather than routing through a Gaussian prior, decomposed into a sequence of adversarial sub-problems for tractability at high resolution. UNSB outperforms CycleGAN and CUT on standard unpaired benchmarks (horse-to-zebra, summer-to-winter) and is architecturally compelling here precisely because it requires no geographic co-registration between domains. It has, however, only been evaluated on natural-image colour-domain shifts, far milder than the shadow-geometry, foreshortening, and occlusion differences between a nadir and a perspective view of the same Mars terrain. Xiao \textit{et al.}~\cite{xiao2025brownian} address the complementary reproducibility problem (stochastic I2I methods yield different outputs on successive runs, unacceptable for scientific visualisation) via a Dual-approx Bridge with separate forward/reverse approximators, achieving deterministic translation with superior SSIM/LPIPS versus stochastic baselines, though again untested on remote sensing imagery.

The methodology (Chapter~3) builds on UNSB as the target Stage~2 architecture, with a CycleGAN baseline (Chapter~4--5) implemented first precisely because, as this section establishes, no reviewed I2I method has been validated against a domain gap this large, making an empirically grounded baseline a necessary first step rather than a formality.

\section{Cross-View Synthesis in Earth Remote Sensing}
\label{sec:crossview}

This cluster was not part of the dissertation's original literature review. It was added after this project's own CycleGAN baseline (Chapter~5, Section~\ref{sec:diagnosis}) failed to translate viewpoint even once its loss-weight and training-shortcut problems were fixed, which prompted a search for how comparable Earth remote-sensing work handles the same nadir-to-ground translation direction. Three papers, none targeting Mars, turned out to address exactly this problem, and none of them attempt pure pixel-space translation.

Qian \textit{et al.}~\cite{qian2023sat2density} propose Sat2Density, which learns a 3D density field from satellite-ground image pairs and renders a ground-level panorama by ray-marching through it, rather than mapping satellite pixels to ground pixels directly. Li \textit{et al.}~\cite{li2024sat2scene} propose Sat2Scene, which generates a 3D sparse voxel representation of urban structure from a satellite image via diffusion before synthesising per-point texture in that representation, again inserting an explicit 3D intermediate rather than translating in image space. Li \textit{et al.}~\cite{li2024crossviewdiff} propose CrossViewDiff, which conditions a diffusion model's denoising process on an explicit satellite-scene structure estimate and a separately computed cross-view texture mapping, rather than relying on the diffusion model to infer structure from the satellite image alone.

All three systems share a pattern absent from this dissertation's original three-stage design: an explicit geometric or structural intermediate sits between the nadir/satellite input and the perspective/ground output, and texture synthesis happens only once that intermediate has resolved the viewpoint change. None treats nadir-to-ground translation as a problem a single conditional network can solve end-to-end from image pairs alone. This is the specific finding that motivated the geometry-mediated redesign reported in Chapter~4 and evaluated in Chapter~5, and the reasoning connecting it to that redesign is given in full in Chapter~3, Section~\ref{sec:deviation-rationale}.

\section{Digital Terrain Modelling and Neural Scene Rendering}

Monocular DTM recovery is a necessary precursor to perspective-view synthesis but does not itself produce visual texture. Tao \textit{et al.}~\cite{tao2021madnet} recover pixel-scale topography from single HiRISE images via photoclinometry (MADNet~2.0); Zou \textit{et al.}~\cite{zou2026pfmgan} improve on this with PFMGAN, an explicit solar-geometry and albedo-aware GAN that reduces reconstruction RMSE by over 50\% across five landform types. This is the strongest quantitative evidence in the survey that physically-grounded architectures outperform unconstrained ones, and a direct design precedent for the physics-informed loss terms in the proposed Stage~2 methodology. Cao \textit{et al.}~\cite{cao2024gen3d} map HiRISE orthoimages to dense DEMs via conditional GAN. None of these three produce rendered surface appearance.

Neural rendering methods are architecturally capable of view-angle transformation but require multi-view input imagery of the \emph{same} scene, which does not exist for a not-yet-visited landing site. Giusti \textit{et al.}~\cite{giusti2022marf} demonstrate NeRF converges on Curiosity rover imagery (MaRF); Li \textit{et al.}~\cite{li2025martian} build a complete text-conditioned video-diffusion simulation pipeline (MarsGen) from over 10,000 annotated Curiosity clips, the most operationally complete system reviewed, but entirely dependent on rover imagery that only exists for four previously-visited sites. Mart\'{i}nez-Petersen \textit{et al.}~\cite{martinezpetersen2025hifi} couple NeRF and Gaussian Splatting in a ROS~2 pipeline, establishing that Gaussian Splatting's interactive render rate is the correct representation target for the VR output this dissertation aims at (Stage~3), even though their pipeline also requires rover imagery as input. Paul \textit{et al.}~\cite{paul2025nerfspace}'s survey confirms, as its central negative finding, that synthesising a NeRF-equivalent scene from a single nadir orbital image has not been addressed anywhere in the literature.

\section{Virtual Reality Reconstruction and Generative Reliability}

Catalano and Soccini~\cite{catalano2025inpainting} inpaint void regions in Mars heightmaps via a boundary-conditioned diffusion model, improving RMSE by up to 15\% and LPIPS by up to 81\% over classical interpolation. A completed heightmap carries no colour information, though, so this addresses only the geometric half of the VR terrain problem, leaving texture synthesis explicitly open.

Every generative model surveyed can, in principle, produce plausible surface features with no basis in physical reality: an aesthetic non-issue for natural images, but a safety-relevant failure mode for mission-critical planetary visualisation. Aithal \textit{et al.}~\cite{aithal2024hallucination} identify \emph{mode interpolation} as the mechanism: a diffusion model's smooth score-function approximation assigns non-zero probability to samples interpolated between nearby training modes (e.g.\ two crater morphologies) that never occurred in training data, a structural property of the approximation, not a defect correctable with more training. Rathkopf~\cite{rathkopf2025reliability} reframes this as an epistemological question, arguing from AlphaFold~3 and GenCast case studies that hallucination risk becomes bounded and scientifically acceptable only through theory-guided training (domain physical constraints) combined with confidence-based error screening. Together, these establish that physical constraint is a primary mechanism for bounding hallucination risk, not an optional refinement. This directly motivates the physics-informed loss terms specified in the methodology (Section~3.2.3) rather than treating Stage~2 as an unconstrained image-translation problem.

\section{Comparative Synthesis}

Table~\ref{tab:comparison} consolidates the twenty-three reviewed publications. Most nadir-domain methods (SR and DTM papers) and perspective-domain methods (NeRF papers evaluated on real Mars data) stay on their own side of the nadir/perspective boundary. The cross-view synthesis cluster (Section~\ref{sec:crossview}) is the exception, and the mechanism by which it crosses that boundary is the load-bearing finding for this dissertation: every one of its three methods inserts an explicit geometric or structural intermediate rather than translating pixels directly, which is not how any of the two Mars-agnostic I2I papers (benchmarked only on natural-image colour-domain shifts) or the original three-stage design approach the problem. The methods with the strongest quantitative improvement within a single domain (PFMGAN, over 50\% RMSE reduction) are also those with explicit physical constraints, while unconstrained SR methods show greater susceptibility to the mode-interpolation artefacts Aithal \textit{et al.} describe, the same pattern the cross-view cluster shows at the translation-architecture level. The two systems closest to a complete Mars VR pipeline, MarsGen and the heightmap-inpainting system, each solve exactly one dimension (perspective texture \emph{or} elevation, never both from nadir input). No reviewed Mars-specific system combines orbital nadir input, perspective-view RGB texture output, and physical-plausibility constraints; the cross-view cluster shows this is achievable in principle, on Earth data, via an explicit intermediate, which is the design pattern Chapter~3, Section~\ref{sec:deviation-rationale} adopts.

\begin{table}[H]
\caption{Comparative overview of the reviewed literature against the nadir-to-rover-perspective VR synthesis problem. Task codes: SR = super-resolution; I2I = image-to-image translation; DTM = digital terrain model generation; NeRF/GS = neural radiance field / Gaussian splatting; VR = VR reconstruction; Data = dataset; Theory = theoretical analysis; Review = survey. View codes: N = nadir/orbital; P = perspective/rover; 3D = 3D representation.}
\label{tab:comparison}
\centering
\resizebox{\textwidth}{!}{%
\renewcommand{\arraystretch}{1.3}
\begin{tabular}{@{}lllllc@{}}
\toprule
\textbf{Ref.} & \textbf{Method} & \textbf{Task} & \textbf{View} & \textbf{Key Result} & \textbf{Mars?} \\
\midrule
\cite{tao2021cassis}    & MARSGAN            & SR      & N$\to$N   & ${\approx}3\times$ SR on CaSSIS, qualitative only          & Yes \\
\cite{wu2024rstsrn}     & RSTSRN              & SR      & N$\to$N   & $+0.88$\,dB PSNR vs.\ LapSRN ($\times4$)                   & Yes \\
\cite{lu2025thermal}    & Multi-source Thermal SR & SR  & N$\to$N   & PSNR gain on THEMIS thermal                                & Yes \\
\cite{khanna2024diffusionsat} & DiffusionSat  & SR/Gen  & N$\to$N   & SOTA FID, multi-sensor                                     & No \\
\cite{zhu2025taming}    & TamingDiff-SR       & SR      & N$\to$N   & $\times4$, PSNR + competitive FID                          & No \\
\cite{kim2024unsb}      & UNSB                & I2I     & Any       & Lower FID vs.\ CycleGAN/CUT, unpaired                      & No \\
\cite{xiao2025brownian} & Dual-approx Bridge  & I2I     & Any       & Deterministic I2I, higher SSIM/LPIPS                       & No \\
\cite{qian2023sat2density} & Sat2Density      & I2I/3D  & N$\to$P   & Density-field render, no depth supervision needed          & No \\
\cite{li2024sat2scene}  & Sat2Scene           & I2I/3D  & N$\to$P   & Diffusion on 3D voxels, photorealistic street sequences    & No \\
\cite{li2024crossviewdiff} & CrossViewDiff    & I2I     & N$\to$P   & Structure+texture-conditioned diffusion, N$\to$P specific  & No \\
\cite{tao2021madnet}    & MADNet~2.0          & DTM     & N$\to$3D  & Monocular DTM $\approx$ stereo accuracy                    & Yes \\
\cite{zou2026pfmgan}    & PFMGAN              & DTM     & N$\to$3D  & $>$50\% RMSE reduction, 5 landform types                   & Yes \\
\cite{cao2024gen3d}     & Cond.\ GAN-DTM      & DTM     & N$\to$3D  & Height consistency, 4 Mars regions                         & Yes \\
\cite{osadnik2025mcted} & MCTED dataset       & Data    & N (pairs) & 80,898 CTX--DEM pairs                                      & Yes \\
\cite{serrano2013hirise}& Bayesian rock density & Analysis & N$\to$N & Probabilistic density from CTX                             & Yes \\
\cite{giusti2022marf}   & MaRF                & NeRF    & P$\to$3D  & Novel view synthesis $\approx$ standard NeRF               & Yes \\
\cite{li2025martian}    & MarsGen             & NeRF+VR & P$\to$VR  & Beats terrestrial video models on FID                      & Yes \\
\cite{dealmeida2025neural3d} & Neural height-field & NeRF & N/oblique & Beats MVS on RMSE (simulated descent)                    & Partial \\
\cite{martinezpetersen2025hifi} & NeRF+GS pipeline & NeRF+GS & P$\to$3D & Real-time GS render, ROS~2 integration                   & Partial \\
\cite{paul2025nerfspace}& NeRF-in-space survey & Review  & ---       & Confirms nadir$\to$perspective gap is open                & Yes \\
\cite{catalano2025inpainting} & Diffusion inpainting & VR & N$\to$VR$_H$ & 15\% RMSE / 81\% LPIPS improvement                     & Yes \\
\cite{aithal2024hallucination} & Mode interpolation & Theory & --- & Mechanistic account of hallucination                    & No \\
\cite{rathkopf2025reliability} & Hallucination-to-reliability & Theory & --- & Physics-guided training bounds risk                & No \\
\bottomrule
\end{tabular}}
\end{table}

\section{Summary}

The literature establishes that each constituent capability needed for VR-grade Mars landing-site visualisation (orbital super-resolution, unpaired image-to-image translation, terrain reconstruction, neural rendering, and hallucination-aware training) has individually advanced substantially since 2021, but no reviewed method integrates them into a pipeline that starts from nadir orbital data and ends with physically-plausible, rover-perspective visual texture. This gap, and the specific finding that physical constraints (not perceptual loss alone) are what separates the strongest-performing methods in each cluster from the weakest, directly motivate the three-stage, physics-constrained architecture formalised in Chapter~3.
